\section{从辅政到接收：权力怎样经过官署、军镇与县廷}

\subsection{辅政不是皇帝身边的一张椅子}

曹叡死后留下幼帝，曹爽与司马懿同受辅政。若只把这一安排读作两个人的性格
冲突，便会错过它为何能牵动淮南、关中和辽东。辅政者所争的是能够让诏令离开
宫廷的接口：宿卫由谁调动，尚书文书由谁经手，近畿军队先听谁的命令，地方
将领又相信哪一方能送来粮饷、任命或赦令。《三国志》魏纪和相关传记可以确定
人事的更替；它们不能让我们把每个官员的选择还原成一场没有沉默的会议。

高平陵之变把这组接口暴露得最清楚。皇帝出城谒陵，洛阳城内外的距离忽然成为
政治条件。司马懿必须使太后名义、城门、宿卫和官署的文书在同一时刻可被使用；
曹爽一方则不得不估量返回、投降或抵抗会如何影响家族和仍在手中的军政位置。
胜利之后的诏令、承诺和指控大多由获胜者保存，不能当作逐时记录；但结果足以
说明，曹魏的实际命令链已经被重排。此后任何边将面对的都不仅是魏帝，而是
谁能让魏帝的名义变成到达军镇的物资与处置。

王凌、毌丘俭、文钦和诸葛诞的行动因此不能串成一条“反司马氏叛乱”的简表。
每一次淮南危机都发生在城池、淮河水道和中枢补给线交错之处。王凌所设想的
继承方案，毌丘俭与文钦所能依赖的军镇，诸葛诞等待的吴援，分别说明不同人
在评估同一问题：洛阳仍是否能够保护他们，或是否已只会要求他们服从。失败
没有证明地方只有一种想法；它证明司马氏在关键时刻比对手更能把粮、兵、诏令
和近畿强制力接到同一处。

\subsection{司马师、司马昭与一份仍用魏名的日常}

司马师废曹芳、立曹髦时，魏朝的官署并未停摆。正因日常任命、文书和地方事务
仍以魏名进行，改立才既能显示辅政者的力量，也能避免立刻承担一切制度断裂的
成本。司马昭接过中军后，淮南围城和关陇准备使军政资源愈发集中。把这种过程
叫作“篡位预演”并非全错，却会使读者误以为禅代已经替每个州郡做完选择。较
准确的说法是，司马氏在魏的制度外壳中不断取得更多可调用入口，外壳仍为日常
秩序提供语言和程序。

曹髦出南阙而被杀，是这个不对称最尖锐的时刻。皇帝并非没有行动意志，也并非
从此可以由一则故事解释所有魏臣；问题在于近侍不足以替代稳定听命的禁军，
宫城之外的道路也已被另一条军政链控制。事后的罪责和礼仪安排，显示谁能决定
如何叙述这场死亡。司马昭仍不立刻受禅，说明名义尚有用途：保留魏帝可让任命、
征发和地方接收继续沿既有程序运作，直到灭蜀带来的军功和资源让公开转换更
容易被安置。

二六五年的魏晋禅代应当读作两种速度。诏册的速度很快，能够改换年号和最高
权属；官署、军镇和家庭所感到的变化却较慢。新朝继承的不只是曹魏的宫门，
还有淮南的军政经验、蜀地尚未安定的接收问题，以及面向孙吴的长期准备。禅代
不是将这些问题清零，而是让西晋有权以新的名义处理它们。

\subsection{孙吴的簿籍如何把江防接到县廷}

走马楼吴简使孙吴的财政不必只作为“江东富庶”或“战时征发”的背景出现。
临湘县的木简留下户、田、租米、吏役、转递和勾核的细部：国家想得到一份粮食
或一次劳力，先要让书手写得下、保管者找得到、基层中介认得出，才可能在江面
或城防上变成可用资源。木简本身并不宏大，正因如此，它让宏大的舰队和宫廷
命令重新有了必须经过的手续。

这些文书不能被强迫回答它们没有提出的问题。临湘并非全吴，残存案卷也不是
连续年度总账；不能由一个县廷算出全国户数、亩数或税率。简上出现的人同样不
等于所有劳动者，未保存的迁徙、依附、逃役与家庭协商并不会因为档案沉默而
消失。它们仍能改变阅读方式：孙权、诸葛恪或孙綝在建业的军政选择，必须在
某些地方被转写成登记、催纳和劳役，才可能支持一段江防。

县廷与江面并非两个世界。沿江运输需要仓口、船只、船工、粮米和知道何时可
行的地方消息；江北有战事时，调度压力会回到县的文书和家庭。孙吴的水路国家
并不是因为长江天然可用，而是不断把这些异质节点勉强接住。卫温、诸葛直远航
的失败尤能说明限度：离开熟悉江岸，风候、疾病、淡水和补给会使朝廷的命令
突然失去可依赖的中介。史书所记的处分能说明朝廷要求结果，不能把夷洲、亶洲
的位置、航程或船员损失补写成确数。

\subsection{钱、役与宫廷的代价}

孙权铸大钱的决定不是一条可脱离生活的经济政策。金属从何而来，工匠如何组织，
市场是否接受，税收与军需怎样结算，都决定一枚钱能否在诏令之外流通。现存记载
不足以重建兑换率或价格趋势，因而不能据此判定孙吴已经陷入或摆脱某种抽象的
“货币危机”。它至少提示朝廷试图在军事、宫廷与地方供给的压力下改变流通手段。

宫室工程和役使争议也应这样读。中枢要营建、守卫和陈列，就会同田间、仓储、
转输和江防争用人力。史料没有连续徭役簿，不能把零星事件化为全吴劳役总量；
但不应因此把成本从叙事里删去。有人被差派、有人须交纳、有人在地方机构中
承担追责，正是国家能力最不平等的部分。吴简不能替建业解释宫廷斗争，宫廷
传记也不能替临湘说明每份簿籍；把两者并置，才可看见同一政权如何在不同尺度
上使资源变得可征、可运和可争。

\subsection{二八〇年以后，接收不是一句“平定江东”}

孙皓出降使建业的最高权属在二八〇年改变。王濬的舟师、杜预的荆州部署和淮南
方向的压力共同迫使吴无法只保住一处渡口；这可由《晋书》《三国志·孙皓传》
和《资治通鉴》的年序互相确认。它们不能让我们把每支舟师的日程、每座城的
态度或每一位吴臣的盘算写成确定事实。降表尤其不能替江东州郡发言。

西晋首先能够接收的是仪式和名义：孙皓被送至洛阳、受封安置，新的朝廷借此把
军事结局纳入自己的礼仪秩序。随后才是较慢的工作。旧吴官吏是否留任，军队
如何遣散或编管，港口的税粮向何处输送，县廷怎样改用新的任命和法律语言，
都要经由地方关系一点点完成。我们没有足以描出每一县接收日期的资料，不应
用统一的“江东归晋”把这种不等速抹平。

对普通家庭而言，政权转换既可能带来新官、新税和新征调，也可能让原有的邻里、
田地和文书方式在一段时间内继续存在。走马楼吴简照亮的是吴时临湘的行政实践，
不能直接证明太康年间全江南的情况；它仍提供一个重要尺度：若一项国家秩序
曾需要如此多登记、收发与勾核，新的秩序就不可能仅凭一份降表在同一天接管。
二八〇年由此既是三国并立的终点，也是西晋开始承担江南日常成本的开端。

\subsection{寿春的堤岸：三次危机怎样改变淮南的日常}

淮河在寿春附近不是一条只供军报标注方向的线。它把城、堤、防营、渡口和通向吴境
的水路放在同一片低地上。谁据寿春，谁便不仅能够在地图上显示一座重镇，也能影响
粮船从哪里停泊、援军向哪一处渡口集中、地方官向谁报送急件。高平陵之后，王凌、
毌丘俭和文钦、诸葛诞相继把这片区域变成反对司马氏的依托，正因为这里原有军政
资源与对吴联络的可能。把三次行动概括为“淮南三叛”，会便利记忆，却会掩去每次
城内外所面对的不同时间表。

嘉平三年王凌案中，密谋、投降及王凌和楚王曹彪的死亡见于《三国志》王凌传和齐王
芳纪。材料能让我们确认有人试图另立宗室，也让人看见司马懿能够把近畿的命令推到
淮南；它没有保存一份所有参与者的名单，更不能证明寿春周围的居民已经在两种政权
之间整齐站队。王凌的方案之所以危险，不在于一纸密约本身，而在于它要求军镇、
使者、道路和可能的宗室名义同时保持秘密而又能行动。消息一旦泄露，原来能保护
计划的河道与城门，也会把投降者的退路压缩为对中枢的服从。

四年后毌丘俭、文钦起兵，问题已经不是简单重复。司马师掌握洛阳的中军，叛军则要
在淮南集合、争取军镇，并判断孙吴是否会越过江淮接应。起兵、失败和文钦南奔可以
由毌丘俭传、文钦传及司马师传交叉确认；兵员数目、吴廷事前知道多少、地方社会
究竟如何反应，并无连续材料可作精确回答。文钦渡向吴境的结果尤其说明，淮南不是
封闭的内地：对一名失败者来说，河流既可能是补给线，也可能是离开魏境的最后通道；
对留在原地的人来说，它却意味着此后更严密的问讯与重新编组。

司马师在镇压后病逝，司马昭接手，并没有让寿春的紧张自动消散。恰恰相反，改由
新辅政者处置军镇，使官员和将领必须重新判断何种命令会被兑现。中枢可通过任免、
赦令和军队展示连续性，但淮南地方所感受到的是更具体的问题：旧将的部曲归谁统领，
仓储应向谁开验，亲属与门生若被卷入前案又能否自保。传记的镜头主要停在将军身上，
我们不应替县廷编造一张完整的恐惧地图；但也不能把中央权力的重新集中写成不触及
地方关系的抽象过程。

甘露二年诸葛诞据寿春向吴求援，把这些已经积累的疑问逼入一场更长的对峙。城内需要
粮，城外的司马昭军需要维持包围，吴军若要救援也必须让船队、人马和消息穿过水网。
《三国志》诸葛诞传、司马昭传、孙亮传和《资治通鉴》都支持叛乱、求援与围城的
基本顺序；它们不足以让我们画出所有营垒、核定城中粮储，或把诸葛诞的每一个决定
归结为个人野心。围城之所以重要，正在于它迫使联盟把承诺变成可抵达的实物，而这
一环最难由口号补足。

孙吴的援助并非一个可以随意写大的箭头。吴廷有牵制魏的理由，但从江东抵近寿春的
队伍仍受航道、接头、时令与魏方封锁限制；而城中每多等一天，军民对粮食、消息和
出路的判断就会改变。寿春最后陷落、诸葛诞被杀、吴援受挫是可以确认的结果。城内
饥困到了什么程度、何时发生何种投降和杀戮、每一支援军损失多少，叙事来源并不
允许把这些痛苦换成确定数字。保留这种限度，并非减轻围城的重量，反而使人看见
史书的胜负语言遮住了多少无法逐户复原的代价。

三次危机的后果不只是司马氏“终于平定淮南”。它改变了中央与地方之间可被相信的
关系。司马昭在寿春取胜后，更能掌握近畿军政与东南压力的接口；对仍在魏名下任官
的人而言，这使公开抵抗的成本骤增。可是把这一结果说成所有人因而真心接受新权力，
同样越过了材料。有人可能在旧岗位上继续做仓储、治安和转运，有人可能换了上司而
保留乡里关系，有人则在战争中失去可见的踪迹。政权把一座城重新纳入命令链，不等
于它同时取得每个人的解释。

寿春还预示了司马氏后来灭蜀、受禅的条件。围城让中枢学会以长时间的补给、隔绝与
接收处理一处军镇，而不只依赖宫廷政变的瞬间。它也留下副作用：要让军镇可靠，
便须加深对将领、部曲和文书的控制；控制越深，地方越可能把安全寄托在少数能够
通向洛阳的人身上。淮河堤岸上的危机结束后，魏国的名号仍在文书中使用，但命令
如何走到城门、仓口与渡口，已经比高平陵以前更清楚地属于司马氏的军政网络。

寿春陷后，最先可见的接收并不是某个抽象的“民心”，而是对军镇人身关系的拆解。诸葛诞的
部曲、随从和原先依附军府的人，不能再只以旧将的名义取得粮食、武器或司法保护。司马昭一方
要处理的，是哪些人应被赦免、哪些人须受审、哪些军士可重新编入、哪些家庭会因亲属在城中而
受到牵连。传记保存了显要人物的处置，未给出一份能够涵盖这些群体的编管册；因此不能替魏廷
宣布一项整齐而迅速的政策。可正是这种史料的不平衡提醒我们，围城结束以后，国家还要把本来
围绕私人军府运转的生活重新接到官署的名义上，而这比攻下一座城慢得多。

淮南人也不只是被动等待接收。渡口和市镇上的人知道哪条水路仍能走，哪些军队已缺粮，亲属在
城内或城外又会怎样影响自家被追问的风险。这类知识极少进入《三国志》的本纪和列传，因为它
们不是为记录地方判断而写；但军队能够迅速围住寿春，本身就依赖这样的消息、船只和地方服务。
不能因为普通人的名字大多没有留下，便把他们写成同一种“叛民”或“顺民”。他们更可能在不同
时段面对不同选择：送粮或藏粮，向谁报告消息，是否跟随一位将领南走，或在新旧命令之间尽量
保住田地与家人。历史不能替他们统一发言，叙事却应承认他们的行动是围城条件的一部分。

寿春的战后安排也令孙吴失去一条原可利用的北方接口。吴廷能够接受文钦等南来者，能够在淮水
附近出兵牵制，却无法只凭一次援军把魏军的包围变成可长期维持的同盟区域。江北的城、仓和
地方关系一旦落入司马昭手中，吴军每次北上都要从更不利的水陆交界重新寻找停泊与接应处。
这不是说孙吴从此没有进攻能力；它说明寿春的失败重排了行动的起点。后来吴廷修筑江防、魏晋
谋划伐蜀和伐吴，都在这个被压缩的东南空间中进行。

对司马氏而言，收束淮南也有不能省略的代价。为了使反叛不再轻易从军镇扩散，中枢须更密集地
把赏罚、任官、监军和物资调配联在一起。这会提高短期的可控制性，却也使许多地方事务更依赖
通向洛阳的少数通道。地方官若首先考虑如何避免被视为与旧将有关，便未必愿意如实报告粮运、
民变或边情；将领若知道一场失败会被写成政治疑点，也会更难公开承认供应不足。史书不能把
这种心理过程逐案核验，但它让我们不应把胜利理解为单纯增加了国家能力。能力的集中同时改变
了什么消息敢于被说、什么关系必须被隐去。

曹髦在二六〇年被杀，之所以能在淮南之后发生，并不因为寿春直接“导致”宫城中的一幕，而是
因为司马昭已在一系列危机中掌握了更稳定的军政入口。皇帝仍有名义和近侍，却缺少能越过城门、
调来可靠军队并把地方资源带进洛阳的链条。淮南的经验告诉司马氏，名义上的魏朝可以继续提供
任免和征发的语言，而实际强制力可由另一套网络集中。两者之间的距离在曹髦事件中变得尖锐，
在二六五年禅代时则被仪式暂时缝合。

把寿春放回这条长线，不是要让三次叛乱成为司马氏必胜的宿命。王凌、毌丘俭、文钦和诸葛诞每
次失败都包含偶然的消息、道路和援军条件；若把它们全改写成权力必然上升，便会抹去当事人在
不完整信息下作出的高风险决定。较可靠的结论更小：连续危机使谁能调动近畿资源、谁能控制
淮水要点、谁能把魏帝名义转为实际命令的问题一次次被公开检验。到寿春陷落时，司马昭比对手
更能通过这些检验；这改变了后来行动的边界，却没有把所有未来结果预先写好。

寿春的堤岸最后留下的是一种制度记忆。防线需要堤、船、粮和人，也需要有人相信下一封文书会
被兑现；叛乱与围城会破坏这种相信，而战后重建它又不能只靠威吓。西晋继承魏的东南部署时，
继承的正是这份昂贵的记忆：如何看守淮河、怎样处置军镇、何时把地方资源集中到一条战线。它
后来帮助新朝面向孙吴，也使新朝必须承担一套已被多次战争磨损的人际与行政成本。
